% ==================== 6.5 本章小结 ====================
\section{本章小结}
\label{sec:ch6_summary}

本章在数字孪生环境下研究了巡视器路径规划方法。主要工作和结论如下：

（1）构建了数字孪生环境下的路径规划框架，采用分层规划架构，包括全局规划层、局部规划层和运动控制层。该框架利用数字孪生中的环境模型、动力学模型和仿真验证能力，实现了更智能的规划决策。

（2）提出了基于注意力机制增强的双深度Q网络（A-D3QN）混合路径规划算法。该算法结合Double DQN、Dueling网络结构和注意力机制，解决了传统DQN的过高估计问题，提高了学习效率和决策质量。

（3）设计了全局-局部协同的混合规划策略。全局规划采用改进的A*算法生成参考路径，局部规划采用A-D3QN进行动态避障和路径修正，两者协同实现了高效、安全的路径生成。

（4）通过实验验证了算法的有效性。在复杂静态场景中，A-D3QN的路径长度比A*缩短1.7\%，成功率提升12个百分点；在动态场景中，成功率达到96\%；路径能耗比A*节省16.2\%。

本章提出的A-D3QN混合路径规划算法为巡视器在月面复杂环境下的自主行走提供了有效的规划方法。
